\section{Applications}
\subsection{Some modally definable classes}
We start by giving some fairly obvious examples of modally definable classes.

Recall that 
\[\Pol(F)\coloneqq \set{P\in\Pol:P\models F}\]

And some standard formulas: 
\begin{definition}$ $\\
    \begin{itemize}
        \item Grz (the Grzegorczyk formula): \[\b(\b(\varphi\to\b\varphi)\to\varphi)\to \varphi\]
        \item Bd$_n$ is the Jankov-Fine formula of the linear order of length $n+1$ 
        \[\text{Bd}_n=\chi(\bullet - \bullet -\cdots -\bullet)\]
        \item $F_n$ (the $n-$fork) is the poset with a single root and $n$ leaves only connected to the root.
    \end{itemize}
\end{definition}
\begin{figure}[H]
    \centering 
    \begin{tikzcd}[ampersand replacement=\&,scale=0.5]
	\bullet \& \bullet \& \bullet \&\& \bullet \& \bullet \& \bullet \& \bullet \& \bullet \\
	\& \bullet \& {F_3} \&\&\&\& \bullet \&\& {F_5}
	\arrow[no head, from=1-1, to=2-2]
	\arrow[no head, from=1-2, to=2-2]
	\arrow[no head, from=1-3, to=2-2]
	\arrow[no head, from=1-5, to=2-7]
	\arrow[no head, from=1-6, to=2-7]
	\arrow[no head, from=1-7, to=2-7]
	\arrow[no head, from=1-8, to=2-7]
	\arrow[no head, from=1-9, to=2-7]
\end{tikzcd}
\caption{A visual representation of the $3-$fork and the $5-$fork}
\end{figure}
\begin{proposition}[Bounded dimension]
    For any given $n\in\N$, the class of polyhedra of dimension
    $\leq n$ is modally definable.
\end{proposition}
\begin{proof}
    Indeed, it is $\Pol(S4.\text{Grz}.\text{Bd}_n)$.
    It was shown in previous works (\cite{DBLP,Dekker}) that a polyhedron has dimension 
    less or equal to $n$ if and only if it validates the bounded depth formula.
\end{proof}

% An interesting (and somewhat surprising) example of a modally definable class is the following:
% \begin{proposition}[Locally $\leq \binom{n}{2}$-holed $1d$]
%     The class of $1d$ polyhedra such that for each point there exist an open neighbourhood
%     has ``fewer than $\binom{n}{2}$ holes''
% \end{proposition}
% \begin{proof}
%     Indeed, $\C=\Pol(S4.\text{Grz}.\text{Bd}_1.\chi(F_{n+1}))$:

%     We know that every $1$ dimensional polyhedron models $S4.\text{Grz}.\text{Bd}_1$; 
%     Let $P\in \C$.
%     Now this means that its face poset contains at most a fork of rank $n$.
%     Let $p$ be the vertex being the root of the fork of highest degree $k\leq n$.
%     Then this means that there are exactly $k$ lines sharing $p$ as a vertex.

%     A hole is obtained by joining at least 2 lines. 
%     Joining $k$ lines to a single point produces $k-1$ holes, which is less than the $\binom{k}{2}$ we would obtain by joining them pairwise.


%     This means that the maximum number of holes that contain $p$ as a point is $\binom{k}{2}$ 
%     (where every pair forms a new hole), meaning the maximum achievable is $\binom{n}{2}$

%     Conversely, if the polyhedron has an open neighbourhood with fewer than $\binom{n}{2}$ holes, 
%     then by reversing the previous argument, its poset contains at most an $n-$fork branching from a vertex.

%     This means that we can take the neighbourhood as the domain for our face up-reduction, and map it to the face poset of 
%     the minial $1-$dimensional polyhedron containing an $n-$fork, by mapping the vertex to the root and each prong to a prong.
% \end{proof}

% \begin{conjecture}
%     Given the proposition above, we have reason to believe that some form of bounded homology on connected components might also be modally definable. Specifically, within the class of polyhedra of bounded dimension $n$, bounding the rank of the $n^{\text{th}}$ homology group might be achievable using Jankov-Fine formulas.
% \end{conjecture}

\subsection{Some classes that are not modally definable}
More interesting are the classes that fail to be modally definable.

Obviously, we have the following.
\begin{proposition}[Connected spaces]
    The class of connected polyhedra is not modally definable.
\end{proposition}
\begin{proof}
    It immediately fails to be closed under disjoint unions.
\end{proof}

\begin{proposition}[Contractible spaces]
    We call a space \textbf{Locally contractible} if every connected component is contractible.

    The class of all (locally) contractible polyhedra is \textbf{not modally
    definable}. 
\end{proposition}
\begin{proof}
    The map illustrated by the drawing below describes a face up-reduction.
  \begin{figure}[H]
    \centering
\begin{tikzpicture}[scale=0.7]
  % Coordinates
  \coordinate (A) at (0.00, 0.00);
  \coordinate (B) at (1.50, 3.00);
  \coordinate (C) at (3.00, 0.00);
  \coordinate (D) at (1.00, 0.50);
  \coordinate (E) at (2.00, 0.50);
  \coordinate (F) at (1.50, 1.75);
  \coordinate (J) at (4.00, 0.00);
  \coordinate (K) at (5.50, 3.00);
  \coordinate (L) at (7.00, 0.00);
  \coordinate (M) at (5.00, 0.50);
  \coordinate (N) at (5.50, 1.75);
  \coordinate (O) at (6.00, 0.50);
  \coordinate (G) at (9.00, 0.50);
  \coordinate (H) at (10.00, 0.50);
  \coordinate (I) at (9.50, 1.75);
  \coordinate (P) at (9.50, 3.00);
  \coordinate (Q) at (8.00, 0.00);
  \coordinate (R) at (11.00, 0.00);
  \coordinate (S) at (12.00, 0.00);
  \coordinate (T) at (15.00, 0.00);
  \coordinate (U) at (13.50, 3.00);

  % Drawing
  \filldraw[pastelred, fill=black, fill opacity=0.3] (A) -- (B) -- (C) -- cycle;
  \draw (D) -- (E) -- (F) -- cycle;
  \draw (F) -- (B);
  \draw (A) -- (D);
  \draw (E) -- (C);
  \filldraw[fill=pastelred, fill opacity=0.3] (A) -- (D) -- (F) -- (B);
  \filldraw[fill=pastelred, fill opacity=0.3] (B) -- (F) -- (E) -- (C);
  \filldraw[fill=pastelred, fill opacity=0.3] (C) -- (E) -- (D) -- (A);
  \draw[pastelred] (J) -- (K) -- (L) -- cycle;
  \draw[pastelred, dashed] (M) -- (N) -- (O) -- cycle;
  \filldraw[pastelred, fill=pastelred, fill opacity=0.3] (M) -- (J) -- (K) -- (N);
  \filldraw[pastelred, fill=pastelred, fill opacity=0.3] (O) -- (L) -- (K) -- (N);
  \filldraw[pastelred, fill=pastelred, fill opacity=0.3] (G) -- (Q) -- (P) -- (I) -- cycle;
  \filldraw[pastelred, fill=pastelred, fill opacity=0.3] (M) -- (J) -- (L) -- (O);
  \filldraw[pastelgreen, fill=pastelgreen, fill opacity=0.3] (R) -- (H) -- (I) -- (P) -- cycle;
  \filldraw[pastelviolet, fill=pastelviolet, fill opacity=0.3] (Q) -- (R) -- (H) -- (G) -- cycle;
  \draw[pastelblue] (I) -- (P);
  \draw[pastelblue] (H) -- (R);
  \draw[pastelblue] (G) -- (Q);
  \draw[pastelred] (S) -- (U);
  \draw[pastelgreen] (T) -- (U);
  \draw[pastelviolet] (S) -- (T);
  \fill[pastelblue] (S) circle (2pt);
  \fill[pastelblue] (T) circle (2pt);
  \fill[pastelblue] (U) circle (2pt);
  \draw[white,dashed] (G) -- (G) -- (H) -- (I) -- cycle;
\end{tikzpicture}
\caption{A simplicial up-reduction $\Delta_2\upred\partial\Delta_2$. In order, $\Delta_2$, the domain of the map, a coloring of the domain and the codomain showing what maps to what.}
\end{figure}
\end{proof}

\begin{proposition}
    The class of convex polyhedra and the class of PL-manifolds are not modally definable.
\end{proposition}
\begin{proof}
    Given a convex polyhedron, we may find a non-convex polyhedron with face poset isomorphic to one of its subdivision, just by moving a vertex 
    to the interior of a maximal face.
    This means that we found a non-convex polyhedron that validates and refutes all the formulas that the original did.

    PL-manifolds (in the sense of \cite{Dekker}\textbf{3}) are even simpler, as they're not closed under disjoint unions.
\end{proof}
